\documentclass{informeutn}
\usepackage[utf8]{inputenc}
\usepackage{array}
\usepackage{geometry}
\usepackage[table]{xcolor}
\usepackage{colortbl}
\usepackage{caption}
\usepackage{graphicx}
\usepackage[label=corner]{karnaugh-map}

\materia{Técnicas digitales I}
\titulo{Trabajo Practico N°1: Algebra de Boole y circuitos combinacionales}
\comision{3R2}
\autores{Ernst Pedro 000000\\ Cortesini Luciano 000000 \\ Prieto Angelo 401012}
\fecha{12-06-2025}

\begin{document}
\maketitle
\tableofcontents

\chapter{Introducción}
  
\chapter{Conversor código BCD a XS3 (exceso 3)}

    \section{Tabla de verdad}

    \begin{table}[h]
        \centering
        \begin{tabular}{cccc|cccc}
            A & B & C & D & $Y_1$ & $Y_2$ & $Y_3$ & $Y_4$ \\
            \hline
            0 & 0 & 0 & 0 & 0 & 0 & 1 & 1 \\
            0 & 0 & 0 & 1 & 0 & 1 & 0 & 0 \\
            0 & 0 & 1 & 0 & 0 & 1 & 0 & 1 \\
            0 & 0 & 1 & 1 & 0 & 1 & 1 & 0 \\
            0 & 1 & 0 & 0 & 0 & 1 & 1 & 1 \\
            0 & 1 & 0 & 1 & 1 & 0 & 0 & 0 \\
            0 & 1 & 1 & 0 & 1 & 0 & 0 & 1 \\
            0 & 1 & 1 & 1 & 1 & 0 & 1 & 0 \\
            1 & 0 & 0 & 0 & 1 & 0 & 1 & 1 \\
            1 & 0 & 0 & 1 & 1 & 1 & 0 & 0 \\
            1 & 0 & 1 & 0 & x & x & x & x \\
            1 & 0 & 1 & 1 & x & x & x & x \\
            1 & 1 & 0 & 0 & x & x & x & x \\
            1 & 1 & 0 & 1 & x & x & x & x \\
            1 & 1 & 1 & 0 & x & x & x & x \\
            1 & 1 & 1 & 1 & x & x & x & x \\
        \end{tabular}
        \caption{Tabla de verdad correspondiente a un conversor de código BCD a XS3.}
        \label{tab:example_label}
    \end{table}

    \section{Funciones lógicas}

    \begin{equation*}
        Y_1 = \sum(5,6,7,8,9)
    \end{equation*}
    \begin{equation*}
        Y_2 = \sum(1,2,3,4,9)
    \end{equation*}
    \begin{equation*}
        Y_1 = \sum(0,3,4,7,8)
    \end{equation*}
    \begin{equation*}
        Y_2 = \sum(0,2,4,6,8)
    \end{equation*}
    \begin{equation*}
        F_x = \sum(10,11,12,13,14,15)
    \end{equation*}

    \section{Minimización}

    \centering
    \begin{tabular}{cc}
      % Primer mapa
      \begin{minipage}{0.45\textwidth}
        \centering
        \begin{karnaugh-map}[4][4][1][$D$][$C$][$B$][$A$]
            \minterms{5,6,7,8,9}
            \terms{10,11,12,13,14,15}{X}
            \implicant{5}{15}
            \implicant{7}{14}
            \implicant{12}{10}
        \end{karnaugh-map}
        \vspace{0.5em}
        $\displaystyle Y_1 = A + B \cdot D + B \cdot C$
      \end{minipage}
      &
      % Segundo mapa
      \begin{minipage}{0.45\textwidth}
        \centering
        \begin{karnaugh-map}[4][4][1][$D$][$C$][$B$][$A$]
            \minterms{1,2,3,4,9}
            \terms{10,11,12,13,14,15}{X}
            \implicant{4}{12}
            \implicantedge{1}{3}{9}{11}
            \implicantedge{3}{2}{11}{10}
        \end{karnaugh-map}
        \vspace{0.5em}
        $\displaystyle Y_2 = B \cdot \overline{C} \cdot \overline{D} + 
                            \overline{B} \cdot D + \overline{B} \cdot C$
      \end{minipage}
      \\[1.5em]
      % Tercer mapa
      \begin{minipage}{0.45\textwidth}
        \centering
        \begin{karnaugh-map}[4][4][1][$D$][$C$][$B$][$A$]
            \minterms{0,3,4,7,8}
            \terms{10,11,12,13,14,15}{X}
            \implicant{0}{8}
            \implicant{3}{11}
        \end{karnaugh-map}
        \vspace{0.5em}
        $\displaystyle Y_3 = \overline{C} \cdot \overline{D} + C \cdot D$
      \end{minipage}
      &
      % Cuarto mapa
      \begin{minipage}{0.45\textwidth}
        \centering
        \begin{karnaugh-map}[4][4][1][$D$][$C$][$B$][$A$]
            \minterms{0,2,4,6,8}
            \terms{10,11,12,13,14,15}{X}
            \implicantedge{0}{8}{2}{10}
        \end{karnaugh-map}
        \vspace{0.5em}
        $\displaystyle Y_4 = \overline{D}$
      \end{minipage}
    \end{tabular}

    \section{Implementación con compuertas}

\chapter{Comparador binario}

\chapter{Conclusión}

\end{document}
